\documentclass{article}

\usepackage{geometry}
\usepackage{makecell}
\usepackage{array}
\usepackage{multicol}
\usepackage{setspace}
\usepackage{changepage}
\usepackage{booktabs}
\usepackage[explicit]{titlesec}
\usepackage{hyperref}
\usepackage{graphicx}
\usepackage{cprotect}
\usepackage{float}
\newcolumntype{?}{!{\vrule width 1pt}}
\newcommand{\paragraphlb}[1]{\paragraph{#1}\mbox{}\\}
\renewcommand{\contentsname}{Inhaltsverzeichnis:}
\renewcommand\theadalign{tl}
\setstretch{1.10}
\setlength{\parindent}{0pt}

\titleformat{\section}
  {\normalfont\Large\bfseries}{\thesection}{1em}{\hyperlink{sec-\thesection}{#1}
\addtocontents{toc}{\protect\hypertarget{sec-\thesection}{}}}
\titleformat{name=\section,numberless}
  {\normalfont\Large\bfseries}{}{0pt}{#1}

\titleformat{\subsection}
  {\normalfont\large\bfseries}{\thesubsection}{1em}{\hyperlink{subsec-\thesubsection}{#1}
\addtocontents{toc}{\protect\hypertarget{subsec-\thesubsection}{}}}
\titleformat{name=\subsection,numberless}
  {\normalfont\large\bfseries}{\thesubsection}{0pt}{#1}

\hypersetup{
    colorlinks,
    citecolor=black,
    filecolor=black,
    linkcolor=black,
    urlcolor=black
}

\geometry{top=12mm, left=1cm, right=2cm}
\title{\vspace{-1cm}Prozessmanagement 1}
\author{Andreas Hofer}

\begin{document}
	\maketitle
	\tableofcontents
	\section{Grundlagen}
  Zuallererst muss man drei Begriffe definieren: Strategie, Taktik und Operatives Handeln.
  \begin{itemize}
    \item{Strategie}
    \begin{itemize}
      \item{Langfristige Planung}
    \end{itemize}
    \item{Taktik}
    \begin{itemize}
      \item{Mittelfristige Planung}
    \end{itemize}
    \item{Operatives Handeln}
    \begin{itemize}
      \item{Kurzfristige Planung}
    \end{itemize}
  \end{itemize}
  Als Beispiel eines Unternehmens würde man bei der Strategie den langfristigen Plan fassen, ein AI-first Unternehmen zu werden. Die Taktik dafür beinhaltet, dass man AI verwendet um einen Kunden automatisiert zur richtigen Region zuzweisen. Das Operative Handeln wäre danach das konkrete Handeln zum Aufbauen und Aufrechterhalten dieses Systems.
  \subsection{Strategie}
  Die Strategie ist das generelle Handeln, welches zu einem Wettbewerbsvorteil führt. Dabei existiert ein Vorteil zwischen entweder der Konkurrenz oder dem Kunden. Für einen Vorteil gegenüber der Kunden, ist oft der Preis relevant, da diese dann eher das eigene Produkt wählen. \\
  \subsubsection{Effektivität vs. Effizienz}
  Hierbei sind auch die Effektivitiät und Effizienz relevant. Effektivität ist der Maßstab wie gut eine Sache vollbracht wird, während die Effizienz betrachtet, wie es mit den geringsten Kosten vollbracht wird. \\
  Effektivitäts- und Effizienzvorteile sind sind die Zeile eine Strategieentwicklung.
  \subsubsection{Vision und Mission}
  Die Strategie benötigt auch eine Vision und eine Mission. Diese Unterscheiden sich darin, für wen sie bestimmt sind. Die Vision ist ein idealerweise kurzes Statement, welches für Mitarbeiter ihre Richtung vorgehen soll und bestimmen, was im allgemeinsten Sinne erreicht werden soll. Walt Disney hat zum Beispiel die Vision 'To make people happy.'. \\
  Die Mission hingegen ist an die Kunden gerichtet und soll kurz mitteilen, was einen erwartet, falls man ihr Produkt ersteht.
  \subsubsection{SMART}
  Um gute Ziele zu entwickeln sollte man diese SMART definieren. Dieses Akronym steht für:
  \begin{itemize}
    \item{Spezifisch}
    \item{Messbar}
    \item{Erreichbar/Attraktiv (Achievable)}
    \item{Angemessen (Reasonable)}
    \item{Terminiert}
  \end{itemize}
  \subsection{Entwicklung}
  Zur Entwicklung der Unternehmensstrategie sollte man diese Maßstäbe verwenden um so seine strategischen Ziele zu definieren und im Endeffekt die Unternehmensstrategie zu definieren. Dabei muss man bei der Umsetzung auch beachten, ob diese Ziele realistisch sind und anhand dessen anpassen. Natürlich können hierbei auch Risiken sichtbar werden, welche man in die Strategie miteinbeziehen muss,
  \subsubsection{Situationsanalyse}
  Zur Entwicklung dieser Unternehmensstrategie sollte man zuerst die Situation analysieren. Dabei sollte man:
  \begin{itemize}
    \item{Zukunftsfaktoren}
    \begin{itemize}
      \item{Wie sieht die momentane und erwartete zukünftige Lage aus?}
    \end{itemize}
    \item{Persönliche Vorteile}
    \begin{itemize}
      \item{Welche Vorteile habe ich als Firma, welche verwendet werden können?}
    \end{itemize}
    \item{[[TODO THIRD ONE]]}
  \end{itemize}
  \subsubsection{Strategieentwicklung}
  Danach sollte man diese Analyse verwenden um seine Strategie konkret zu entwickeln und ein Strategiestatement zu erarbeiten. Dabei muss man beachten, wie man von seine Analyse zu einer Strategie kommt und was benötigt wird um es zu vollbringen.
  \subsubsection{Strategiestatement}
  Ein Strategiestatement hat oft ein festes Schema, um seine Ziele überschaubar weitergeben zu können. \\
  Ein Statement besteht meist aus diesen Teilen:
  \begin{itemize}
    \item{Kern}
    \begin{itemize}
      \item{Wie sieht die momentane Lage aus?}
    \end{itemize}
    \item{Ziele}
    \begin{itemize}
      \item{Wo wollen wir hin?}
    \end{itemize}
    \item{Wirkungsbereich}
    \begin{itemize}
      \item{Was muss geschehen um dort hin zu kommen?}
    \end{itemize}
    \item{Kundennutzen}
    \begin{itemize}
      \item{Was ist der Nutzen gegenüber der Kunden, wenn es erreicht wird?}
    \end{itemize}
    \item{Unser Vorteil}
    \begin{itemize}
      \item{Was ist der Vorteil den das Unternehmen aus der Strategie gewinnen soll?}
    \end{itemize}
  \end{itemize}
  \subsubsection{Strategieumsetzung}
  Wenn man die Strategie erarbeitet hat, mus man diese nun umsetzen. Dazu benötigt man eine Strategy Map welches die Zusammenhänge der verschiedenen Ziele definieren soll. Danach muss man die Messgrößen festlegen um Messen zu können, wie der Erfolg definiert wird. Als letztes müssen die Zielwerte bestimmt werden um definieren zu können, wann man sein Ziel erreicht hat. Aus diesen Überlegungen sollen dann Projekte entstehen, welche die Strategie umsetzen. Die Umsetzung geschieht jedoch selten genau so wie erwartet weshalb man immer dafür planen muss, dass es Abweichungen gibt und diese Ausgleichen.
  \subsection{Grundstrategien nach Porter}
  Die Grundstrategien nach Porter gelten in der Regel für jedes Geschäftsfeld und unterteilen Unternehmen in spezifische Kategorien. Dabei kombiniert man jeweils den Strategischen Vorteil und die Ausrichtung des Unternehmens. Bei dem Strategischen Vorteil unterscheidet man zwischen dem Kundennutzen oder dem Kostenvorteil. Bei der Ausrichtung unterscheidet man zwischen einer Ausrichtung der gesamten Branche oder nur einem einzelnen Segment.
  \begin{tabular}{| l | l | l | l | l |}
    \toprule
    \multicolumn{4}{Strategischer Vorteil} \\ \hline
    &&Hervorragender Kundennutzen & Kostenvorteile \\ \hline
    
    
    \bottomrule
  \end{tabular}
  Bei Fokus auf die gesamte Branche und dem Vorteil eines Kundennutzens spricht man einer Differenzierung. Bei Fokus auf den Kostenvorteil spricht man von einer Kostenführerschaft. Wenn man stattdessen auf ein einzelnes Segment ausgerichtet ist, spricht man von einer Schwerpunktbildung unabhängig ob man Kundennutzen oder Kostenvorteile bildet. In der Regel ist es eine schlechte Idee alle dieser vier Segmente gleichzeitig behandeln zu wollen.
  \subsubsection{Differenzierung}
  Bei der Differenzierung versucht das Unternehmen sich gegenüber Kunden so zu differenzieren, dass dieses gewählt wird. Das kann durch eine überdurchschnittliche Qualität, Kundenservice, Markenimage oder anderes gegeben werden. Dabei dürfen Kosten jedoch nicht ignoriert werden, sind jedoch nicht das Hauptziel. In der Regel haben Unternehmen in diesem Segment hohe Forschungs- und Entwicklungskosten, da sie ihre Marktstellung bewahren wollen, oder qualifiziertes Personal bzw. hohe Rohstoffqualität.
  \subsubsection{Kostenführerschaft}
  Die Kehrseite davon ist die Kostenführerschaft, wo man einen Wettbewerbsvorteil durch einen Kostenvorteil erreichen will. Dabei muss man nicht auch gleichzeitig die Preisführerschaft innehaben, diese ist jedoch meist Vorraussetzung. Viele der Kosten sind meist relativ statisch, weshalb es schwierig sein kann Kosten zu senken. Eine Strategie kann es sein durch großes Volumen die Anschaffungskosten zu verringern.
  \subsubsection{Schwerpunktsbildung}
  Bei der Schwerpunktsbildung fokussiert man sich auf bestimmte Käufergruppen, Produktsegmente, oder geografische Märkte. Als Unternehmen sollte man dadurch sich auf das eigene Marktsegment fokussieren um diesen Marktvorteil zu erhalten.
  \subsubsection{Stuck in the Middle}
  Wenn man keine der anderen Strategien abdeckt, ist man in der Regel \textit{Stuck in the Middle}. Dies ist der schlechteste Ort zu sein, da man oft hohe Kosten durch ein breites Angebot hat, ohne aber dessen Kostenvorteile zu genießen. Ein Unternehmen sollte stets dessen Strategie überprüfen, ob man sich nicht durch einen veränderten Markt sich nun in der Mitte befindet.

  \subsection{Five Forces nach Porter}
  Nach Porter gibt es fünf Kräfte, welche eine Situation auf dem Markt beeinflussen. Diese sind:
  \begin{itemize}
    \item{Vorhandener Wettbewerb}
    \item{Lieferanten}
    \item{Kunden}
    \item{Potentielle Wettbewerber}
    \item{Substitution der Produkte}
  \end{itemize}
  \subsubsection{Vorhandener Wettbewerb}
  Wenn man der einzige Hersteller eines spezifischen Produkts ist, erfährt man keine Rivalität. Da das jedoch selten der Fall ist hat man oft Mitbewerber. Je mehr es davon gibt, umso größer ist die Rivalität.
  \subsubsection{Verhandlungsposition des Kunden}
  Kunden können gegenüber eines Unternehmens auch eine Verhandlungsposition besitzen. Das ist gegeben wenn:
  \begin{itemize}
    \item{Kunden über die Marktusituation informiert sind}
    \item{Diese einen großen Teil des Umsatzes des Händlers haben}
    \item{Der Kunde das Produkt selbst herstellt (Backward Integration)}
  \end{itemize}
  \subsubsection{Verhandlungsposition des Lieferanten}
  Lieferanten können auch Druck auf ein Unternehmen aufbauen wenn:
  \begin{itemize}
    \item{Geringe Anzahl an Lieferanten für ein spezielles Produkt}
    \item{Das Unternehmen unbedeutend ist}
    \item{Der Lieferant das Produkt selbst herstellt (Forward Integration)}
  \end{itemize}
  \subsubsection{Ersatzprodukte}
  Wenn ein Produkt durch eine alternative Methode einfacher, günstiger oder besser bereitgestellt werden kann, kann eine Substitution entstehen. Ein Beispiel dafür ist der Videoverleih, welcher durch das Aufkommen von Online Streamingdiensten nahezu komplett verdrängt wurde, da diese das praktisch gleiche Produkt schneller und meist günstiger zur Verfügung stellen konnten.
  \subsection{SWOT-Analyse}
  Das SWOT in SWOT-Analyse steht für: Stärken (Strengths), Schwächen (Weaknesses), Chancen (Opportunities) und Risiken (Risks). \\
  Eine SWOT-Analyse wird durchgeführt indem man zuerst die IST Situation untersucht und danach das auf zukünftige Entwicklung erweitert. Die Istsituation besteht aus einer Unternehmensanalyse sowie einer Mitbewerberanalyse. Dabei wird untersucht wie man als Unternehmen positioniert ist und wie Rivalen im gleichen oder berührenden Sektor auf einen einwirken. Bei der zukünftigen Entwicklung muss man eine Umweltanalyse durchführen und herausfinden, wie die Trends aussehen und was das für das Unternehmen bedeutend könnte. Es gibt verschiedenste Versionen wie man
  \subsubsection{Unternehmensanalyse}
  \section{Prozessmanagement}
  Prozessmanagement beschreibt die Betrachtung der Aktivitäten eines Unternehmens. Die Umwelt des Prozessmanagements ist der PDCA-Zyklus (Plan, Do, Check, Act):
  \begin{enumerate}
    \item{Zuerst muss analysiert werden}
    \item{Danach müssen Definitionen geschaffen und ein Modell erstellt werden}
    \item{Dieses muss danach implementiert werden}
    \item{Diese Implementierung muss danach durchgeführt und gesteuert werden}
    \item{Zu Ende muss die Implementierung überprüft und optimiert werden}
  \end{enumerate}
  \subsection{Reifegradmodell}
  Man kann jedoch nicht sofort von null auf hundert beim Prozessmanagement gehen, weshalb das Reifegradmodell modellieren soll, wo man als Unternehmen beim Prozessmanagement gerade steht. Dabei gibt es fünf Stufen:
  \begin{enumerate}
    \item{Rudimentäre und zufällige Abläufe}
    \begin{itemize}
      \item{Aubläufe haben keinen festen Faden und werden nach Bedarf abgearbeitet (Man fragt einen Kollegen ob er schnell etwas für einen erledigen kann)}
    \end{itemize}
    \item{Wiederholte Aubläufe (Prozeduren)}
    \begin{itemize}
      \item{Es besteht eine gewisse Systematik bei den Abläufen (Man fragt stets den selben Kollegen)}
    \end{itemize}
    \item{Strategiegerecht definierte Geschäftsprozesse}
    \begin{itemize}
      \item{Die erste Stufe bei der man wirklich von Prozessmanagement redet. Hierbei besteht nun eine Methodik bei den Abläufen (Man fragt den Kollegen der dafür verantwortlich ist)}
    \end{itemize}
    \item{Gemessene Geschäftsprozesse}
    \begin{itemize}
      \item{Der Fokus besteht nicht mehr auf den Abläufen selbst sondern deren Effizienz}
    \end{itemize}
    \item{Optimierte Prozesse}
    \begin{itemize}
      \item{Prozesse werden kontinuierlich gemessen, überprüft und verbessert.}
    \end{itemize}
  \end{enumerate}
  \subsection{Auslöser und Erwartungen}
  Doch wieso benötigt man überhaupt Prozessmanagement? Abläufe innerhalb eines Unternehmens zu definieren und optimieren kann zu einer Steigerung der Produktivität und Qualität, aber auch neuer Innovation und einer besseren Transparenz des betrieblichen Geschehens führen. Typische Erfolge einer solchen Implementierung sind:
  \subsection{Erfolgsrelevanz}
  Der Erfolg zur Einführung eines Prozesses muss natürlich messbar sein. Verbesserungen sollten in genau definierten Erfolgsgrößen oder zugrundeliegenden Prozessfähigkeiten sichtbar gemacht werden. Wenn der Prozess zum Beispiel nicht linear mit dem Endprodukt verbunden werden kann, muss man das relativieren.
  \subsection{Personenunabhängigkeit}
  Die Mitarbeiter einer jeden Firma sind meist dessen wertvollste Resource. Prozesse können zwar Abläufe optimieren, wenn diese jedoch nicht von Mitarbeitern angenommen oder implementiert werden, ist dessen Effektivität eingeschränkt. Aus diesem Grund sollten Prozesse idealerweise unabhängig von den einzelnen Personen implementierbar sein.
  \subsection{Konsistenz}
  Varianten in einem Prozess (Zum Beispiel Varianten bei Automobilen) können in einem Prozess zu großem Aufwand führen. Wenn man zum Beispiel ein Unternehmen mit vielen Varianten anbietet, kann die Konsistenz des Angebots leiden. Aus diesem Grund sollte man stets überprüfen wie viele Varianten eines Produkts man wirklich anbieten muss.
  \subsection{Anpassungsfähigkeit}
  Mit der Zeit kann sich die Effektivität und Effizienz von Prozessen verändern. Diese Veränderung kann auch schleichend geschehen, weshalb man immer die Prozesse analysieren und optimieren muss.
  \subsection{Offenheit}
  Da Prozessmanagement versucht einen Ablauf in eine geordnete Bahn zu lenken und Innovation das Ziel hat neue Pfade zu beschreiten, ist es eventuell manchmal schwer zu vereinen. Dabei muss man jedoch beachten, dass nicht jede Innovation stets eine radikale Änderung bedeutend muss. Es kann auch inkrementelle Innovationen geben, welche neues beschreiten aber nicht sofort eine radikale Änderung bedeuten. Natürlich kann es auch radikale Änderungen geben, dabei muss man jedoch wahrscheinlich neue Prozesse definieren.
  \subsection{Begriffserklärung}
  \subsection{Beobachtungen aus der Praxis}
  Trotz eines guten Optimierungsvorhabens, scheitern diese oft bei der Umsetzung, wodurch die Theorie an der Praxis scheitert.
  \subsection{Strategie- und Prozessmanagement}
  \subsection{Change Management}
  Change Management beschreibt den Prozess
  \section{Makromodellierung}
  Die Basiseinheit des Makromodells ist der Geschäftsprozess. Dieser wird durch einen einzelnen Geschäftsfall aktiviert. Die Aktivitäten des Prozesses aktivieren sich im Idealfall in einer kontrollierten Reihenfolge. Als Beispiel einer Werbeanzeige wird zuerst die Anzeige in Auftrag gegeben, was den Prozess beginnt (Geschäftsfall), die Prozesse reihen sich danach aneinander was im Endeffekt zur paginierten Seite führt. Jeder Schritt in diesem Prozess erzeugt dabei einen Mehrwert für das Endprodukt und wenn alles richtig verläuft, sollte das Produkt im Endeffekt einen größeren Mehrwert bieten, als die Kosten für den Prozess betragen. 
  \subsection{Prozesseigner}
  Um die End-to-End Verantwortung zu garantieren, sollte in jedem dieser Schritte die gleiche Abteilung verantwortlich sein. Diese Rolle nennt man Prozesseigner, welcher von Anfang des Projekts bis zu dessen Ende für die Umsetzung verantwortlich ist.
  \subsection{Wertschaffend und Wertdefinierend}
  Man unterscheidet zwischen vier Arten von Prozessen:
  \begin{itemize}
    \item{Wertschaffender Prozess}
    \begin{itemize}
      \item{Direkte Gewinnung und Abwicklung von Aufträgen}
    \end{itemize}
    \item{Wertdefinierender Prozess}
    \begin{itemize}
      \item{Marketing sowie Gestaltung der Prozesse und deren Entwicklung}
    \end{itemize}
    \item{Managementprozesse}
    \begin{itemize}
      \item{Definierung der Strategie oder Finanzplanung und Budgetierung sowie Controlling und HR}
    \end{itemize}
    \item{Supportprozess}
    \begin{itemize}
      \item{Administration, Buchhaltung und IT}
    \end{itemize}
  \end{itemize}
  Diese vier Prozesse kann man in zwei Dimensionen beschreiben: Einerseits wie oft diese wiederholt werden, andererseits wie weit dessen Entscheidungsspielraum geht. Im unteren linken Eck befinden sich dabei die Wertschaffenden Prozesse, da genau definiert ist, was in dem Prozess passieren muss und diese werden auch oft wiederholt, wobei in der Regel nur der Prozess selbst involviert ist. \\
  Supportprozesse befinden sich in der gleichen Ecke, weisen jedoch in der Regel einen höheren Spielraum auf und müssen idealerweise nicht so oft wiederholt werden.
  \subsection{Erarbeitung}
  Um ein Makrodesign zu erarbeiten, durchgeht man in der Regel einen Prozess aus fünf Schritten:
  \begin{enumerate}
    \item{Marktleistung und Architektur}
    \begin{itemize}
      \item{Darstellung der Kundenbedürfnisse sowie aller Leistungen}
    \end{itemize}
    \item{Integraler Leistungsprozess}
    \item{Kaskadierung}
    \item{Segmentierung}
    \item{Konsolidierung}
  \end{enumerate}
  \subsubsection{Marktleistung}
  Das Marktleistungsbündel beschreibt den Prozess für den Kunden um alle Leistungen anhand der Kundenbedürfnisse zu befriedigen. Man unterscheidet hierbei zwischen der straffen und der losen Kopplung. Bei der \textbf{straffen Kopplung} ist die Vorleistung maßgeschneidert, weshalb eine Leistung nicht mehr oder nur schlecht erbracht werden kann, wenn ein anderer Prozess sich verändert. Dem gegenüber steht die \textbf{lose Kopplung} wobei die Leistungserbringungen unabhängig voneinander sind weshalb die einzelnen Leistungen stark variieren können und die Leistung trotzdem erbracht werden können. Hauptaugenmerk des Marktleistungsbündels ist, ein Verständnis für die Kundenbedürfnisse zu gewinnen und diese optimal befriedigen zu können. \\
  Idealerweise durchgeht die Interaktion mit dem Kunden einem Zyklus bei der stetig eine Austauschbeziehung besteht um die Leistung zu gewährleisten.
  \subsubsection{Integraler Leistungsprozess}
  \subsubsection{Kaskadierung}
  Bei der Kaskadierung löst ein Geschäftsfall einen Subprozess aus, welcher den ersten Prozess unterstützt. Zum Beispiel kann die Bestellung in einem Webshop eine Reihe an Subprozessen ausführen. Eine Bestellung führt dazu, dass die Erstellung des Produkts angefordert wird, was wiederum den Versand anfordert, welches im Endeffekt an den Kunden geliefert wird.
  \paragraphlb{Bildung}
  Eine Kaskadierung wird durch vier Schritte gebildet:
  \begin{enumerate}
    \item{Neutrale Prozesskette}
    \begin{itemize}
      \item{Alle Prozessschritte werden aneinandergereiht}
    \end{itemize}
    \item{Prozesskettenplan}
    \begin{itemize}
      \item{Die Schritte werden in Aktivitäten gespalten}
    \end{itemize}
    \item{Bildung der durchgängigen Prozessverantwortung}
    \begin{itemize}
      \item{Es wird festgestellt, wer für jede Aktivität verantwortlich ist}
    \end{itemize}
    \item{Bildung der Kaskade}
    \begin{itemize}
      \item{Anhand der Verantwortungen können gewisse Aktivitäten so als Kaskade ausgelagert werden}
    \end{itemize}
  \end{enumerate}
  \paragraphlb{Echte Kaskadierung}
  Die echte Kaskadierung ist nur gegeben, wenn zwischen Anfang und Ende des Subprozesses keine Interaktion zwischen den beiden Prozessen passiert, da die gesamten Feinheiten der Erwartungen bereits geklärt wurden.
  \paragraphlb{Unechte Prozesskaskade}
  Eine unechte Prozesskaskade ist das Gegenteil der echten Kaskade, wo der Subprozess immer wieder mit dem Hauptprozess interagiert. Es kann Vorteile haben eine unechte Prozesskaskade aufzubauen, wenn zum Beispiel regelmäßige Rücksprache nötig ist, jedoch ist es in vielen Fällen ratsam diesen Subprozess stattdessen in den Hauptprozess zu integrieren da die Effizienz bei einer unechten Kaskade in der Regel signifikant sind. \\
  Man muss jedoch bedenken, dass eine echte Kaskadierung oft unrealistisch ist, da ein Auftrag sich stets wegen Änderungen der Kundenwünsche ändern kann, was eine Interaktion mit dem Subprozess nötig macht.
  \paragraphlb{Mehr-Augen-Prinzip}
  Eine Abwandlung dieser Kaskadierung ist eine inbegriffene Qualitätssicherung. So kann der Subprozess nur angefordert werden, wenn eine separate Überprüfung zum Beispiel der Kreditwürdigkeit nötig ist.
  \subsubsection{Segmentierung}
  Die Segmentierung versucht für Prozessvarianten spezifische Vorgänge einzuführen um so abhängig von den Anforderungen effizienter agieren zu können. Das kann relevant werden, wenn zwei Varianten in großen Teilen gleich sind, sich jedoch in einigen Teilen unterschieden. Wenn diese Prozesse segmentiert sind, muss so nicht jedes Mal überlegt werden wie es abgeändert werden kann
  \subsubsection{Konsolidierung}
  Die Konsolidierung überprüft vergangene Aufspaltungsprozesse um diese zu kombinieren, um die Effizienz zu steigern. Das kann nötig sein, wenn zwei diskrete Prozesse große Synergien besitzen, und so eine Vereinigung den Overhead reduziert.





















  
\end{document}